\documentclass{report}
\usepackage{amsfonts, amsmath}
\newcommand\R{{\mathbb R}}
\newcommand\Q{{\mathbb Q}}
\newcommand\N{{\mathbb N}}
\pagestyle{empty}
\begin{document}
\large
\centerline{\Huge Analisi Matematica II}
\centerline{\large Corso di Laurea di Informatica}
\renewcommand{\thefootnote}{ } 
\footnote{\sl Avete 2.30 ore di tempo. Ogni esercizio
vale nove punti. La sufficienza si ottiene con un punteggio $\geq$
18. Solo le risposte chiaramente giustificate saranno prese in
considerazione.}
\renewcommand{\thefootnote}{\arabic{footnote}} 
\vskip.1cm
\centerline{\Large Appello 9-06-2003}
\vskip1cm
\begin{enumerate}
\item Si determini il raggio di convergenza della serie di potenze
\[
\sum_{n=0}^\infty\left\{[e^{3n}+e^n]^{\frac 13}-e^n\right\}x^n.
\]
\item Si scriva la serie di Taylor per la funzione $f(x)=(1+2x^2)^{-1}$. 
\item Per ogni $a>0$ sia $x_a(t)$ la soluzione del problema di Cauchy
\[
x'=x-ax^2
\]
\[
x(0)=1.
\]
Si verifichi che $x_a(t)$ \`e ben definito per tutti i $t>0$. Infine,
si calcoli
\[
\lim_{t\to\infty}x_a(t).
\]
\item Sia $f:[-\pi,\pi]\to\R$ definita da
\[
f(x)=-1\quad \text{per }x\leq 0
\]
\[
f(x)=  1\quad \text{per }x>0.
\]
Si calcoli la serie di Fourier di $f$ e si usi il risultato per
calcolare la somma dei numeri dispari a segni alterni, ovvero della serie
\[
\sum_{n=0}^\infty\frac{(-1)^n}{2n+1}.
\]
\end{enumerate}
\end{document}
